% Figure: magnitude response of a second-order notch (band-reject) filter,
% |H| = |1 - (f/f0)^2| / sqrt( (1 - (f/f0)^2)^2 + (f/(Q f0))^2 ), drawn at
% two quality factors. The response is unity in both pass bands and dips to
% a deep null at f0; the higher-Q curve gives a narrower, sharper notch.
% Frequency normalised to f/f0 on a log scale.
\begin{figure}[htbp]
\centering
\begin{tikzpicture}
\begin{axis}[
  width=12cm, height=7.2cm,
  xmode=log,
  xlabel={normalised frequency $f/f_0$}, ylabel={magnitude (dB)},
  xmin=0.1, xmax=10, ymin=-45, ymax=5,
  xtick={0.1,0.2,0.5,1,2,5,10},
  xticklabels={0.1,0.2,0.5,1,2,5,10},
  ytick={0,-10,-20,-30,-40},
  axis lines=left,
  tick label style={font=\scriptsize},
  label style={font=\small},
  legend pos=south east,
  legend style={font=\scriptsize},
  clip=false,
]
% Q = 1 (broad notch). Guard the null by clamping the argument away from 0
% is unnecessary here because the sample grid avoids exactly f/f0 = 1; the
% log-spaced 300 samples over 0.1..10 do not land on 1 exactly.
\addplot[engblue, very thick, domain=0.1:10, samples=300]
  {20*log10( abs(1 - x^2) / sqrt( (1 - x^2)^2 + (x/1)^2 ) )};
\addlegendentry{$Q=1$ (broad notch)}
% Q = 8 (sharp notch)
\addplot[engorange, very thick, domain=0.1:10, samples=300]
  {20*log10( abs(1 - x^2) / sqrt( (1 - x^2)^2 + (x/8)^2 ) )};
\addlegendentry{$Q=8$ (sharp notch)}
% pass-band reference line
\addplot[enggray, dashed, domain=0.1:10, samples=2] {0};
% mark the notch centre
\draw[enggray, dotted] (axis cs:1,-45) -- (axis cs:1,0);
\node[enggray, font=\scriptsize, anchor=south] at (axis cs:1,0.4) {$f_0$};
\end{axis}
\end{tikzpicture}
\caption[Notch (band-reject) filter magnitude at two quality factors]{Magnitude
response of a second-order notch filter at two quality factors. Both curves pass
unchanged well below and well above the notch centre and dip toward a deep null
at the centre frequency, where the numerator vanishes. The higher-quality-factor
curve rejects a narrow band around the centre and returns to the pass band
quickly, the response wanted for removing a single interfering tone such as
mains hum while disturbing nearby signal frequencies as little as possible; the
lower-quality-factor curve carves a broad trough that attenuates a wider band at
the cost of touching more of the wanted spectrum. The depth of the null is set
by how precisely the bridge elements balance, so a real notch reaches thirty to
forty decibels rather than the infinite rejection the ideal formula promises,
and a tunable notch trims one element to place the null exactly on the offending
frequency.}
\label{fig:vol04ch09:notch-response}
\end{figure}
